\documentclass[a4paper,12pt]{report}

% Paket-paket yang diperlukan
\usepackage[bahasa]{babel}
\usepackage{graphicx}
\usepackage{hyperref}
\usepackage{enumitem}
\usepackage{listings}
\usepackage{xcolor}
\usepackage{tcolorbox}
\usepackage{courier}
\usepackage{fancyhdr}
\usepackage{indentfirst}


% Package TikZ untuk diagram
\usepackage{tikz}
\usepackage{pgf}
\usetikzlibrary{mindmap}
\usetikzlibrary{arrows.meta}
\usetikzlibrary{positioning}
\usetikzlibrary{shapes,arrows}
\usetikzlibrary{er}
\usetikzlibrary{calc}

% Library tambahan untuk diagram kompleks
\usetikzlibrary{trees}
\usetikzlibrary{shadows}
\usetikzlibrary{backgrounds}
\usetikzlibrary{decorations}

% Atur penomoran section agar tidak terpengaruh chapter
% Ini akan memungkinkan penomoran section dimulai dari 1 di setiap chapter
\counterwithout{section}{chapter}

% Konfigurasi hyperref
\hypersetup{
    colorlinks=true,
    linkcolor=black,
    filecolor=magenta,      
    urlcolor=cyan,
    pdftitle={Modul Praktikum Sistem Manajemen Basis Data},
    pdfpagemode=FullScreen,
}

% Konfigurasi untuk kode program
\definecolor{codegreen}{rgb}{0,0.6,0}
\definecolor{codegray}{rgb}{0.5,0.5,0.5}
\definecolor{codepurple}{rgb}{0.58,0,0.82}
\definecolor{backcolour}{rgb}{0.95,0.95,0.92}

\lstdefinestyle{mystyle}{
    backgroundcolor=\color{backcolour},
    commentstyle=\color{codegreen},
    keywordstyle=\color{magenta},
    numberstyle=\tiny\color{codegray},
    stringstyle=\color{codepurple},
    basicstyle=\ttfamily\small,
    breakatwhitespace=false,
    breaklines=true,
    captionpos=b,
    keepspaces=true,
    numbers=left,
    numbersep=5pt,
    showspaces=false,
    showstringspaces=false,
    showtabs=false,
    tabsize=2
}

\lstset{style=mystyle}

% Konfigurasi header dan footer
\pagestyle{fancy}
\fancyhf{}
\fancyhead[L]{}
\fancyhead[R]{Departemen Teknik Komputer ITS}
\fancyfoot[C]{\thepage}
\renewcommand{\headrulewidth}{0.4pt}
\renewcommand{\footrulewidth}{0.4pt}

\begin{document}

% Modifikasi penomoran ToC untuk tidak menampilkan 0 di section
\renewcommand{\thesection}{\arabic{section}}

% Input Cover
% Informasi dokumen untuk cover
\title{\Huge \textbf{Modul Praktikum\\Basis Data: ACID dan High Concurrency\\dengan PostgreSQL}}
\author{Departemen Teknik Komputer\\Institut Teknologi Sepuluh Nopember}
\date{2025}

% Halaman cover
\begin{titlepage}
    \centering
    \vspace*{1cm}
    {\includegraphics[width=0.3\textwidth]{logo-tekkom.png}\par}
    \vspace{1.5cm}
    {\huge\bfseries Modul Praktikum\\Sistem Manajemen Basis Data\par}
    \vspace{1cm}
    {\Large  ACID dan High Concurrency PostgreSQL\par}
    \vspace{2cm}
    {\Large\itshape Departemen Teknik Komputer\\Institut Teknologi Sepuluh Nopember\\Surabaya\par}
    \vfill
    {\large Penulis: Arta Kusuma Hernanda\par}
    \vspace{0.5cm}
    % Footer
    {\large 2025\par}
\end{titlepage}

% Input Kata Pengantar
% Kata Pengantar
\chapter*{Kata Pengantar}
\addcontentsline{toc}{chapter}{Kata Pengantar}

Puji syukur kami panjatkan kepada Tuhan Yang Maha Esa atas segala rahmat dan karunia-Nya sehingga Modul Praktikum Basis Data: ACID dan High Concurrency dengan PostgreSQL ini dapat tersusun hingga selesai. Modul ini disusun untuk memenuhi kebutuhan mahasiswa dalam pembelajaran mata kuliah Basis Data di Departemen Teknik Komputer, Institut Teknologi Sepuluh Nopember Surabaya.

Modul praktikum ini bertujuan untuk memberikan pemahaman mendalam tentang konsep ACID (Atomicity, Consistency, Isolation, Durability) dan High Concurrency pada database PostgreSQL melalui empat skenario praktis. Melalui modul ini, mahasiswa diharapkan dapat memahami dan mempraktikkan berbagai teknik optimasi database seperti indexing, transactions, partitioning, serta strategi backup dan recovery.

Kami menyadari bahwa modul praktikum ini masih jauh dari kesempurnaan. Oleh karena itu, kritik dan saran yang membangun sangat kami harapkan demi perbaikan modul ini di masa yang akan datang.

Akhir kata, kami berharap modul praktikum ini dapat bermanfaat bagi mahasiswa dalam mengembangkan kemampuan dalam pengelolaan basis data dan dapat menjadi bekal yang berguna di dunia kerja nantinya.

\vspace{1cm}

\begin{flushright}
Surabaya, April 2025\\

Tim Penyusun
\end{flushright}

% Daftar Isi
\tableofcontents
\newpage

% Input Persiapan Praktikum (konten README.md)

% Gunakan chapter* agar tidak bernomor
\chapter*{Persiapan Praktikum}
\addcontentsline{toc}{chapter}{Persiapan Praktikum}

\renewcommand{\thesection}{\arabic{section}}

\section{Deskripsi Praktikum}

Praktikum ini bertujuan untuk memahami konsep ACID (Atomicity, Consistency, Isolation, Durability) dan High Concurrency pada database PostgreSQL melalui empat skenario praktis.

Praktikum ini terdiri dari empat modul yang saling berkaitan:

\begin{enumerate}
    \item \textbf{Indexing}: Mempelajari cara mengoptimalkan performa query dengan index dan mengukur dampaknya pada konkurensi tinggi.

    \item \textbf{Transaction}: Mempelajari konsep ACID, tingkat isolasi transaksi, dan menangani konkurensi tinggi dengan benar.

    \item \textbf{Partitioning}: Mempelajari cara membagi tabel besar menjadi bagian yang lebih kecil dan dampaknya terhadap performa.

    \item \textbf{Backup dan Recovery}: Mempelajari strategi backup dan recovery untuk menjamin ketersediaan dan integritas data.
\end{enumerate}

%\section{Struktur Repositori}

\section{Persiapan Environment}

\subsection{Menggunakan Docker}

\begin{enumerate}
    \item Install Docker dan Docker Compose di komputer Anda
    \item Unduh file praktikum yang sudah disediakan oleh asisten praktikum
    \item Jalankan docker-compose file untuk membuat container dengan perintah:
\end{enumerate}

\begin{lstlisting}[language=bash]
docker-compose up -d
\end{lstlisting}

\begin{enumerate}
    \setcounter{enumi}{3}
    \item Generate data dummy:
\end{enumerate}

\begin{lstlisting}[language=bash]
cd scripts
pip install faker
python generate_data.py
\end{lstlisting}

\begin{enumerate}
    \setcounter{enumi}{4}
    \item Import data ke PostgreSQL:
\end{enumerate}

\begin{lstlisting}[language=bash]
cd scripts
bash import_data.sh
\end{lstlisting}

\begin{enumerate}
    \setcounter{enumi}{5}
    \item Akses pgAdmin di browser:
    \begin{itemize}
        \item URL: \url{http://localhost:5050}
        \item Email: admin@example.com
        \item Password: p4ssw0rd
    \end{itemize}

    \item Tambahkan server di pgAdmin:
    \begin{itemize}
        \item Name: PostgreSQL Praktikum
        \item Host: postgres
        \item Port: 5432
        \item Username: praktikan
        \item Password: p4ssw0rd
    \end{itemize}
\end{enumerate}

\section{Petunjuk Penggunaan Modul}

Setiap modul praktikum memiliki:
\begin{itemize}
    \item PDF Modul
    \item File README.md dengan penjelasan konsep dan langkah-langkah
    \item File SQL dengan query untuk latihan
    \item Tugas 
\end{itemize}

\section{Panduan Pengerjaan}

\begin{enumerate}
    \item Praktikum dikerjakan dalam kelompok 
    \item Setiap kelompok mengerjakan semua modul praktikum
    \item Pelaksanaan praktikum dilakukan oleh asisten praktikum dengan pengawasan koordinator dosen praktikum
    \item Setiap kelompok membuat laporan akhir berisi:
    \begin{itemize}
        \item Jawaban tugas dari setiap modul
        \item Analisis dan kesimpulan
    \end{itemize}
\end{enumerate}

\section{Struktur Folder Praktikum}

Folder praktikum memiliki struktur sebagai berikut:

\begin{lstlisting}[language=bash,basicstyle=\ttfamily\small]
	postgres-praktikum/
	|-- docker-compose.yml           
	|-- init/                         
	|   |-- 01-schema.sql            
	|   `-- 02-functions.sql        
	|-- data/                       
	|   |-- products.csv             
	|   |-- customers.csv           
	|   |-- orders.csv            
	|   `-- order_items.csv           
	|-- scripts/                    
	|   |-- generate_data.py          
	|   |-- import_data.sh           
	|   |-- benchmark.sh             
	|   `-- README.md               
	`-- praktikum/                   
	|		|-- praktikum1_indexing/      
	|		|-- praktikum2_transaction/ 
	|		|-- praktikum3_partitioning/ 
	`--	|-- praktikum4_backup/    
\end{lstlisting}

\section{Dataset Praktikum}

Dataset yang digunakan adalah data e-commerce dummy yang berisi:
\begin{itemize}
    \item 10.000 produk
    \item 5.000 pelanggan
    \item 20.000 pesanan
    \item 50.000 item pesanan
\end{itemize}

\section{Referensi}

\begin{itemize}
    \item \href{https://www.postgresql.org/docs/}{PostgreSQL Documentation}
    \item \href{https://www.postgresql.org/docs/current/mvcc.html}{PostgreSQL Concurrency}
    \item \href{https://www.postgresql.org/docs/current/indexes.html}{PostgreSQL Indexing}
    \item \href{https://www.postgresql.org/docs/current/ddl-partitioning.html}{PostgreSQL Partitioning}
    \item \href{https://www.postgresql.org/docs/current/backup.html}{PostgreSQL Backup and Restore}
\end{itemize}

% Modul P1
\setcounter{section}{0}
\chapter{Implementasi Indexing pada PostgreSQL}
\setcounter{section}{0}
\section{Pendahuluan}
Index pada database adalah struktur data khusus yang mempercepat operasi pencarian dan pengambilan data. Analoginya seperti indeks pada buku yang membantu kita menemukan konten tertentu dengan cepat tanpa perlu membaca seluruh buku. Dalam PostgreSQL, index disimpan dalam struktur khusus yang mengoptimalkan pencarian berdasarkan kolom tertentu.

Implementasi indexing yang tepat merupakan salah satu strategi utama untuk meningkatkan performa database, terutama pada sistem dengan data besar dan tingkat konkurensi tinggi.

\subsection{Tujuan Praktikum}
Dilaksanakannya praktikum indexing, praktikan diharapkan mampu:
\begin{enumerate}
    \item Memahami konsep dasar indexing pada database PostgreSQL
    \item Mempelajari tipe-tipe index di PostgreSQL
    \item Mengimplementasikan index yang tepat untuk sebuah skenario
    \item Menganalisis dampak index terhadap performa query konkuren
\end{enumerate}

\section{Konsep Dasar Indexing}

\subsection{Apa itu Index?}
Index pada database bekerja mirip dengan indeks buku: membantu sistem menemukan data dengan cepat tanpa memeriksa seluruh tabel (table scan). Index menyimpan pointer ke baris data dalam tabel berdasarkan nilai kolom yang diindeks.

\subsection{Bagaimana Index Bekerja}

Ketika query dijalankan, PostgreSQL menganalisis apakah menggunakan index akan lebih efisien daripada melakukan table scan. Jika menggunakan index lebih efisien, PostgreSQL akan:
\begin{enumerate}
    \item Mencari data di index berdasarkan kondisi WHERE
    \item Menemukan referensi (pointer) ke baris data yang relevan
    \item Mengambil data dari tabel menggunakan pointer tersebut
\end{enumerate}

\subsection{Kapan Menggunakan Index}
Index sangat berguna pada:
\begin{enumerate}
    \item Kolom yang sering digunakan dalam klausa WHERE
    \item Kolom yang sering digunakan untuk JOIN antar tabel
    \item Kolom yang sering digunakan dalam klausa ORDER BY atau GROUP BY
    \item Tabel berukuran besar dengan query yang hanya mengambil sebagian kecil data
\end{enumerate}

\subsection{Kapan Tidak Menggunakan Index}
Index tidak selalu bermanfaat pada:
\begin{enumerate}
    \item Tabel kecil yang lebih efisien dilakukan sequential scan
    \item Kolom dengan kardinalitas rendah (nilai yang berbeda sedikit)
    \item Kolom yang jarang digunakan dalam query
    \item Tabel yang sering diupdate/dihapus (overhead pemeliharaan index)
\end{enumerate}

\subsection{Dampak Index pada Performa}
Index memberikan dampak pada:
\begin{enumerate}
    \item \textbf{SELECT}: Mempercepat query dengan mengurangi jumlah baris yang perlu diperiksa
    \item \textbf{INSERT}: Memerlukan overhead tambahan untuk memperbarui index
    \item \textbf{UPDATE}: Memerlukan overhead untuk memperbarui index jika kolom yang diindex berubah
    \item \textbf{DELETE}: Memerlukan overhead untuk memperbarui index
\end{enumerate}

\section{Jenis-jenis Index di PostgreSQL}

PostgreSQL menyediakan beberapa jenis index yang dapat dipilih sesuai dengan kebutuhan:

\subsection{B-tree Index}
\begin{enumerate}
    \item Index default di PostgreSQL
    \item Efisien untuk perbandingan dengan operator equality (=) dan range $(<, >, <=, >=, BETWEEN)$
    \item Mendukung urutan untuk ORDER BY
    \item Cocok untuk kebanyakan kasus penggunaan
\end{enumerate}

\subsection{Hash Index}
\begin{enumerate}
    \item Optimal hanya untuk operator equality (=)
    \item Lebih cepat dari B-tree untuk operasi equality sederhana
    \item Tidak mendukung range queries atau sorting
    \item Dirancang untuk tabel hash di memori
\end{enumerate}

\subsection{GiST Index (Generalized Search Tree)}
\begin{enumerate}
    \item Index untuk data geometri, text-search, dan data kompleks lainnya
    \item Fleksibel, dapat digunakan untuk data custom
    \item Mendukung nearest-neighbor searches
    \item Digunakan untuk full-text search, data geografis, dll
\end{enumerate}

\subsection{GIN Index (Generalized Inverted Index)}
\begin{enumerate}
    \item Dirancang untuk nilai yang memiliki multiple components (arrays, jsonb, text-search)
    \item Efisien untuk pencarian yang membutuhkan matching multiple values
    \item Cocok untuk kolom yang menyimpan data semi-structured
    \item Lebih lambat untuk operasi insert dibanding GiST
\end{enumerate}

\subsection{BRIN Index (Block Range INdex)}
\begin{enumerate}
    \item Dirancang untuk tabel sangat besar dengan data yang terurut secara natural
    \item Sangat kecil dan efisien untuk kolom seperti timestamp, ID sequential, dll
    \item Kinerja query lebih rendah dari B-tree tapi overhead penyimpanan jauh lebih kecil
\end{enumerate}

\subsection{SP-GiST Index (Space-Partitioned GiST)}
\begin{enumerate}
    \item Untuk data yang bisa dipartisi secara non-overlapping
    \item Baik untuk data hierarchical seperti rentang IP, geo-data
    \item Mendukung nearest-neighbor searches
\end{enumerate}

\begin{figure}[h]
	\centering
	\begin{tabular}{|l|p{10cm}|}
		\hline
		\textbf{Jenis Index} & \textbf{Karakteristik dan Penggunaan} \\
		\hline
		B-tree & Index default PostgreSQL. Efisien untuk operasi perbandingan (=, <, >, BETWEEN) dan mendukung pengurutan (ORDER BY). \\
		\hline
		Hash & Khusus untuk operasi equality (=). Lebih cepat dari B-tree untuk lookup sederhana, tidak mendukung range. \\
		\hline
		GiST/GIN & Untuk data kompleks seperti full-text search, data spasial, array, dan JSON. GIN lebih lambat untuk insert tapi lebih cepat untuk search. \\
		\hline
		BRIN & Untuk tabel besar dengan data terurut secara natural (seperti timestamp). Ukuran kecil, overhead rendah. \\
		\hline
	\end{tabular}
	\caption{Jenis-jenis Index di PostgreSQL dan Penggunaannya}
\end{figure}

\subsection{Benchmark pada Konkurensi Tinggi}

Pada situasi konkurensi tinggi, performa index dapat dipengaruhi oleh beberapa faktor:

\begin{itemize}
    \item \textbf{Contention}: Kompetisi antar koneksi untuk mengakses data yang sama
    \item \textbf{Lock contention}: Waktu tunggu karena lock pada baris atau tabel
    \item \textbf{Index bloat}: Overhead karena index yang jarang di-maintenance
\end{itemize}

\section{Best Practices Implementasi Indexing}

\begin{enumerate}
    \item \textbf{Index kolom yang sering digunakan dalam WHERE, JOIN, ORDER BY}
    \item \textbf{Hindari over-indexing}: Terlalu banyak index berdampak negatif pada performa INSERT/UPDATE/DELETE
    \item \textbf{Gunakan composite index} untuk query dengan multiple conditions
    \item \textbf{Perhatikan urutan kolom} pada composite index (most selective first)
    \item \textbf{Pertimbangkan partial index} untuk subset data yang sering diakses
    \item \textbf{Gunakan EXPLAIN ANALYZE} untuk validasi penggunaan index
    \item \textbf{Jalankan VACUUM ANALYZE} secara berkala untuk memperbarui statistik
    \item \textbf{Monitor ukuran dan penggunaan index} menggunakan pg\_stat\_*
    \item \textbf{Lakukan REINDEX} pada index yang terfragmentasi
    \item \textbf{Evaluasi trade-off} antara kecepatan query vs overhead pemeliharaan
\end{enumerate}

\section{Tahapan Praktikum}

\subsection{Persiapan}
\begin{enumerate}
	\item Pastikan container Docker PostgreSQL sudah berjalan
	\item Gunakan pgAdmin atau PSQL untuk mengakses database
	\item Pastikan data sudah diimpor menggunakan script yang disediakan
\end{enumerate}

\subsection{Mengeksplorasi Execution Plan}
Pertama, Anda akan mempelajari cara melihat execution plan query menggunakan EXPLAIN dan EXPLAIN ANALYZE:

\begin{lstlisting}[language=SQL]
	-- Menampilkan execution plan
	EXPLAIN SELECT * FROM products WHERE category = 'Elektronik';
	
	-- Menampilkan execution plan beserta runtime
	EXPLAIN ANALYZE SELECT * FROM products WHERE category = 'Elektronik';
\end{lstlisting}

\subsection{Mengukur Performa Query Tanpa Index}
\begin{enumerate}
	\item Jalankan query berikut tanpa index dan catat waktu eksekusinya:
	
	\begin{lstlisting}[language=SQL]
		-- Query 1: Filter berdasarkan kategori
		SELECT COUNT(*) FROM products WHERE category = 'Elektronik';
		
		-- Query 2: Filter berdasarkan range harga
		SELECT * FROM products WHERE price BETWEEN 1000000 AND 5000000;
		
		-- Query 3: Filter berdasarkan tanggal pesanan
		SELECT o.order_id, o.customer_id, o.order_date 
		FROM orders o 
		WHERE o.order_date BETWEEN '2023-06-01' AND '2023-06-30';
		
		-- Query 4: Join tanpa index
		SELECT c.first_name, c.last_name, o.order_id, o.order_date, o.total_amount
		FROM customers c 
		JOIN orders o ON c.customer_id = o.customer_id 
		WHERE c.city = 'Jakarta';
		
		-- Query 5: Aggregate yang berat
		SELECT p.category, COUNT(*) as total_products, AVG(p.price) as avg_price
		FROM products p
		GROUP BY p.category
		ORDER BY total_products DESC;
	\end{lstlisting}
	
	\item Lakukan benchmark konkuren tanpa index:
	
	\begin{lstlisting}[language=bash]
		# Di terminal
		bash query_before.sql 20 100 10
	\end{lstlisting}
\end{enumerate}

\subsection{Implementasi Index yang Tepat}
Sekarang, tambahkan index yang sesuai untuk setiap query:

\begin{lstlisting}[language=SQL]
	-- Hapus index yang sudah ada (untuk tujuan praktikum)
	DROP INDEX IF EXISTS idx_products_category;
	DROP INDEX IF EXISTS idx_products_price;
	DROP INDEX IF EXISTS idx_orders_date;
	DROP INDEX IF EXISTS idx_customers_city;
	
	-- Index 1: B-Tree index untuk kategori produk
	CREATE INDEX idx_products_category ON products(category);
	
	-- Index 2: B-Tree index untuk range harga
	CREATE INDEX idx_products_price ON products(price);
	
	-- Index 3: B-Tree index untuk tanggal pesanan
	CREATE INDEX idx_orders_date ON orders(order_date);
	
	-- Index 4: Index untuk city pada tabel customers
	CREATE INDEX idx_customers_city ON customers(city);
	
	-- Index 5: Composite index untuk JOIN operation
	CREATE INDEX idx_orders_customer_id ON orders(customer_id);
\end{lstlisting}

\subsection{Mengukur Performa Setelah Indexing}
\begin{enumerate}
	\item Jalankan kembali query yang sama dan bandingkan waktu eksekusinya:
	
	\begin{lstlisting} [language=SQL]
		-- Jalankan semua query yang sama seperti sebelumnya dan bandingkan execution plan
		EXPLAIN ANALYZE SELECT COUNT(*) FROM products WHERE category = 'Elektronik';
		-- dan seterusnya untuk query lainnya
	\end{lstlisting}
	
	\item Lakukan benchmark konkuren dengan index:
	
	\begin{lstlisting} [language=bash]
		# Di terminal
		cd scripts
		bash benchmark.sh ../praktikum/praktikum1_indexing/query_after.sql 20 100 10
	\end{lstlisting}
\end{enumerate}

\subsection{Mempelajari Tipe Index Lain}
PostgreSQL mendukung beberapa tipe index. Cobalah implementasi berikut:

\begin{lstlisting} [language=SQL]
	-- Index Partial untuk produk dengan stok rendah
	CREATE INDEX idx_products_low_stock ON products(product_id, stock_quantity) 
	WHERE stock_quantity < 10;
	
	-- Index menggunakan ekspresi untuk pencarian case-insensitive
	CREATE INDEX idx_products_name_lower ON products(LOWER(name));
	
	-- BRIN Index untuk data berurutan (seperti timestamp)
	CREATE INDEX idx_orders_date_brin ON orders USING BRIN(order_date);
	
	-- GIN Index untuk pencarian full-text (jika menggunakan PostgreSQL >= 9.6)
	CREATE INDEX idx_products_description_gin ON products 
	USING GIN(to_tsvector('english', description));
\end{lstlisting}

\section{Referensi}

\begin{itemize}
    \item \href{https://www.postgresql.org/docs/current/indexes.html}{PostgreSQL Documentation: Indexes}
    \item \href{https://www.postgresql.org/docs/current/using-explain.html}{PostgreSQL Documentation: EXPLAIN}
    \item \href{https://www.postgresql.org/docs/current/performance-tips.html}{PostgreSQL Documentation: Performance Tips}
    \item \href{https://www.postgresql.org/docs/current/indexes-types.html}{PostgreSQL Documentation: Index Types}
    \item \href{https://www.postgresql.org/docs/current/pgbench.html}{PostgreSQL Documentation: pgbench}
\end{itemize}

% Modul P2
\setcounter{section}{0}
\chapter{Implementasi Transaction dan ACID pada PostgreSQL}
\setcounter{section}{0}
\section{Pendahuluan}
Transaksi (transaction) merupakan salah satu konsep fundamental dalam sistem basis data relasional yang menjamin integritas dan reliabilitas data. Dalam PostgreSQL, implementasi transaksi yang tepat dan pemahaman ACID properties sangat penting untuk mengembangkan aplikasi yang handal, terutama dalam lingkungan dengan konkurensi tinggi.

Transaksi adalah unit kerja yang mengubah data dari satu state ke state lainnya. Semua perubahan dalam transaksi harus berhasil secara lengkap, atau tidak ada yang dilakukan sama sekali (all-or-nothing). Konsep ini menjamin bahwa database tetap dalam keadaan konsisten meskipun terjadi kegagalan sistem.

\subsection{Tujuan Praktikum}
Dilaksanakannya praktikum transaction dan ACID, praktikan diharapkan mampu:
\begin{enumerate}
    \item Memahami konsep ACID dalam transaksi database
    \item Mengimplementasikan transaksi pada PostgreSQL
    \item Mempelajari tingkat isolasi transaksi
    \item Menganalisis fenomena konkurensi seperti dirty read, non-repeatable read, dan phantom read
    \item Mengimplementasikan strategi penanganan konkurensi tinggi
\end{enumerate}

\section{Konsep ACID dalam Transaksi Database}

ACID adalah akronim yang merepresentasikan empat properti penting dari transaksi database yang menjamin reliabilitas data meskipun terjadi kegagalan sistem, error, atau konkurensi akses.

\subsection{Atomicity (Atomisitas)}
\begin{enumerate}
    \item Transaksi harus dijalankan sepenuhnya atau dibatalkan sepenuhnya
    \item Tidak ada hasil parsial; jika satu operasi gagal, semua operasi dibatalkan
    \item Implementasi menggunakan mekanisme COMMIT dan ROLLBACK
    \item Contoh: Transfer uang antar rekening harus mengurangi saldo satu rekening dan menambah rekening lain, atau tidak melakukan perubahan sama sekali
\end{enumerate}

\subsection{Consistency (Konsistensi)}
\begin{enumerate}
    \item Transaksi harus membawa database dari satu state valid ke state valid lainnya
    \item Semua aturan integritas data (constraints, cascade, triggers) harus tetap terpenuhi
    \item Jumlah total saldo sebelum dan sesudah transaksi harus tetap sama
    \item Melindungi aplikasi dari data yang tidak konsisten
\end{enumerate}

\subsection{Isolation (Isolasi)}
\begin{enumerate}
    \item Transaksi harus berjalan seolah-olah tidak ada transaksi lain yang berjalan secara bersamaan
    \item Mencegah transaksi membaca data "kotor" dari transaksi lain yang belum commit
    \item PostgreSQL mengimplementasikan multi-version concurrency control (MVCC)
    \item Terdapat beberapa tingkat isolasi yang bisa dipilih sesuai kebutuhan
\end{enumerate}

\subsection{Durability (Durabilitas)}
\begin{enumerate}
    \item Perubahan yang sudah di-commit harus tetap tersimpan, bahkan jika terjadi kegagalan sistem
    \item PostgreSQL menggunakan Write-Ahead Logging (WAL) untuk menjamin durabilitas
    \item Data yang sudah di-commit harus bertahan dari crash, power failure, dll
    \item Memungkinkan pemulihan data setelah kegagalan sistem
\end{enumerate}

\section{Tingkat Isolasi Transaksi}

PostgreSQL mendukung beberapa tingkat isolasi sesuai dengan standar SQL. Setiap tingkat isolasi menawarkan keseimbangan antara konsistensi dan performa.

\subsection{READ UNCOMMITTED}
\begin{enumerate}
    \item PostgreSQL tidak benar-benar mengimplementasikan tingkat ini
    \item Dalam PostgreSQL, READ UNCOMMITTED sama dengan READ COMMITTED
    \item Tidak mungkin membaca perubahan yang belum di-commit (dirty read)
\end{enumerate}

\subsection{READ COMMITTED (Default)}
\begin{enumerate}
    \item Tingkat isolasi default di PostgreSQL
    \item Transaksi hanya dapat membaca data yang sudah di-commit
    \item Mencegah dirty read, tetapi memungkinkan non-repeatable read dan phantom read
    \item Membaca yang sama dalam satu transaksi dapat mengembalikan hasil berbeda jika transaksi lain melakukan commit
\end{enumerate}

\subsection{REPEATABLE READ}
\begin{enumerate}
    \item Transaksi hanya melihat data yang sudah di-commit sebelum transaksi dimulai
    \item Mencegah dirty read dan non-repeatable read
    \item Melindungi dari perubahan data yang dibaca, tetapi masih memungkinkan phantom read
    \item Semua query dalam satu transaksi melihat snapshot database yang sama
\end{enumerate}

\subsection{SERIALIZABLE}
\begin{enumerate}
    \item Tingkat isolasi tertinggi
    \item Mencegah semua anomali konkurensi (dirty read, non-repeatable read, phantom read)
    \item Transaksi berjalan seolah-olah dieksekusi secara serial (satu per satu)
    \item Dapat menyebabkan serialization error yang memerlukan retry
    \item Memberikan konsistensi tertinggi dengan trade-off performa
\end{enumerate}

\begin{figure}[h]
	\centering
	\begin{tabular}{|l|c|c|c|}
		\hline
		\textbf{Tingkat Isolasi} & \textbf{Dirty Read} & \textbf{Non-repeatable Read} & \textbf{Phantom Read} \\
		\hline
		READ UNCOMMITTED* & Tidak & Ya & Ya \\
		\hline
		READ COMMITTED & Tidak & Ya & Ya \\
		\hline
		REPEATABLE READ & Tidak & Tidak & Ya** \\
		\hline
		SERIALIZABLE & Tidak & Tidak & Tidak \\
		\hline
	\end{tabular}
	\caption{Tingkat Isolasi dan Anomali Konkurensi di PostgreSQL\\
	*Sama dengan READ COMMITTED di PostgreSQL\\
	**Di PostgreSQL, REPEATABLE READ juga mencegah phantom read, yang melebihi standar SQL}
\end{figure}

\subsection{Anomali Konkurensi}

Berikut adalah penjelasan lebih detail mengenai anomali yang dapat terjadi:

\begin{enumerate}
    \item \textbf{Dirty Read}: Membaca data yang belum di-commit oleh transaksi lain
    \item \textbf{Non-repeatable Read}: Transaksi membaca data dua kali, tetapi antara dua pembacaan tersebut, data diubah dan di-commit oleh transaksi lain
    \item \textbf{Phantom Read}: Transaksi mengeksekusi query yang mengembalikan himpunan baris, tetapi transaksi lain insert/delete baris yang memenuhi kondisi query
\end{enumerate}

\section{Strategi Penanganan Konkurensi}

Untuk menangani konkurensi dengan baik, PostgreSQL menyediakan beberapa strategi:

\subsection{Pessimistic Concurrency Control}
\begin{enumerate}
    \item Menggunakan locking untuk mencegah konflik
    \item Row-level locking dengan FOR UPDATE, FOR SHARE, dll
    \item Table-level locking dengan LOCK TABLE
    \item Mencegah konflik sebelum terjadi
    \item Dapat menyebabkan deadlock
\end{enumerate}

\subsection{Optimistic Concurrency Control}
\begin{enumerate}
    \item Mengasumsikan konflik jarang terjadi
    \item Tidak mengunci data, tetapi memeriksa apakah data berubah sebelum commit
    \item Biasanya menggunakan version number atau timestamp
    \item Cocok untuk lingkungan dengan read-heavy workload
    \item Memerlukan retry jika terjadi konflik
\end{enumerate}

\subsection{Multi-Version Concurrency Control (MVCC)}
\begin{enumerate}
    \item PostgreSQL menggunakan MVCC untuk implementasi isolasi
    \item Setiap transaksi melihat snapshot database pada titik waktu tertentu
    \item Memungkinkan reader tidak mengunci writer dan sebaliknya
    \item Meningkatkan konkurensi karena read tidak menghalangi write
    \item Trade-off adalah overhead penyimpanan untuk multiple versions
\end{enumerate}

\section{Tahapan Praktikum}

\subsection{Persiapan}
\begin{enumerate}
	\item Pastikan container Docker PostgreSQL sudah berjalan
	\item Gunakan pgAdmin atau PSQL untuk mengakses database
	\item Pastikan data sudah diimpor menggunakan script yang disediakan
\end{enumerate}

\subsection{Implementasi Transaksi Dasar}
Mari kita mulai dengan mempelajari cara mengimplementasikan transaksi dasar di PostgreSQL:

\begin{lstlisting}[language=SQL]
	-- Transaksi dasar
	BEGIN;
	UPDATE products SET stock_quantity = stock_quantity - 5 WHERE product_id = 1;
	INSERT INTO orders (customer_id, status) VALUES (1, 'pending');
	COMMIT;
	
	-- Transaksi dengan ROLLBACK
	BEGIN;
	UPDATE products SET stock_quantity = stock_quantity - 100 WHERE product_id = 1;
	-- Oops, stok tidak cukup!
	ROLLBACK;
\end{lstlisting}

\subsection{Demonstrasi Tingkat Isolasi Transaksi}
Selanjutnya, kita akan mendemonstrasikan perbedaan antara berbagai tingkat isolasi:

\subsubsection{READ COMMITTED (Default)}
Buka dua terminal PostgreSQL secara bersamaan dan jalankan:

Terminal 1:
\begin{lstlisting}[language=SQL]
	BEGIN;
	SET TRANSACTION ISOLATION LEVEL READ COMMITTED;
	UPDATE bank_accounts SET balance = balance - 1000000 WHERE account_id = 1;
	-- Tunggu beberapa saat sebelum commit
	SELECT pg_sleep(30);
	COMMIT;
\end{lstlisting}

Terminal 2:
\begin{lstlisting}[language=SQL]
	BEGIN;
	SET TRANSACTION ISOLATION LEVEL READ COMMITTED;
	-- Ini akan menampilkan nilai sebelum update di Terminal 1
	SELECT * FROM bank_accounts WHERE account_id = 1;
	-- Tunggu Terminal 1 commit
	SELECT pg_sleep(35);
	-- Ini akan menampilkan nilai setelah update di Terminal 1
	SELECT * FROM bank_accounts WHERE account_id = 1;
	COMMIT;
\end{lstlisting}

\subsubsection{REPEATABLE READ}
Buka dua terminal PostgreSQL dan jalankan:

Terminal 1:
\begin{lstlisting}[language=SQL]
	BEGIN;
	SET TRANSACTION ISOLATION LEVEL REPEATABLE READ;
	SELECT * FROM bank_accounts WHERE account_id = 1;
	-- Tunggu Terminal 2 melakukan update dan commit
	SELECT pg_sleep(30);
	-- Ini masih akan menampilkan nilai asli meskipun Terminal 2 sudah commit
	SELECT * FROM bank_accounts WHERE account_id = 1;
	COMMIT;
\end{lstlisting}

Terminal 2:
\begin{lstlisting}[language=SQL]
	BEGIN;
	UPDATE bank_accounts SET balance = balance - 1000000 WHERE account_id = 1;
	COMMIT;
\end{lstlisting}

\subsubsection{SERIALIZABLE}
Buka dua terminal PostgreSQL dan jalankan:

Terminal 1:
\begin{lstlisting}[language=SQL]
	BEGIN;
	SET TRANSACTION ISOLATION LEVEL SERIALIZABLE;
	SELECT SUM(balance) FROM bank_accounts WHERE customer_id = 1;
	-- Tunggu Terminal 2 menambahkan account baru
	SELECT pg_sleep(30);
	-- Ini tidak akan melihat account baru karena isolasi serializable
	SELECT SUM(balance) FROM bank_accounts WHERE customer_id = 1;
	COMMIT;
\end{lstlisting}

Terminal 2:
\begin{lstlisting}[language=SQL]
	BEGIN;
	INSERT INTO bank_accounts (customer_id, balance)
	VALUES (1, 5000000);
	COMMIT;
\end{lstlisting}

\subsection{Mengatasi Kondisi Konkurensi Tinggi}

\subsubsection{Menggunakan Lock (Pessimistic Concurrency Control)}
\begin{lstlisting}[language=SQL]
	-- Row-level lock
	BEGIN;
	SELECT * FROM products WHERE product_id = 1 FOR UPDATE;
	-- Sekarang row ini terkunci hingga transaksi selesai
	UPDATE products SET stock_quantity = stock_quantity - 5 WHERE product_id = 1;
	COMMIT;
	
	-- Lock pada tingkat tabel
	BEGIN;
	LOCK TABLE products IN EXCLUSIVE MODE;
	-- Operasi pada tabel products
	COMMIT;
\end{lstlisting}

\subsubsection{Optimistic Concurrency Control}
\begin{lstlisting}[language=SQL]
	-- Menggunakan last_updated sebagai version control
	BEGIN;
	SELECT product_id, stock_quantity, last_updated
	FROM products WHERE product_id = 1;
	
	-- Asumsikan kita mendapatkan last_updated = '2023-07-15 10:00:00'
	UPDATE products
	SET stock_quantity = stock_quantity - 5,
	    last_updated = CURRENT_TIMESTAMP
	WHERE product_id = 1 AND last_updated = '2023-07-15 10:00:00';
	
	-- Jika tidak ada baris yang terupdate, berarti data sudah berubah
	-- dan kita harus mengulang transaksi
	COMMIT;
\end{lstlisting}

\subsection{Implementasi Function Transfer Uang}
\begin{lstlisting}[language=SQL]
	-- Function sudah dibuat di 02-functions.sql
	-- Cobalah dari dua terminal berbeda
	BEGIN;
	SELECT * FROM bank_accounts WHERE account_id IN (1, 2);
	SELECT transfer_money(1, 2, 500000);
	SELECT * FROM bank_accounts WHERE account_id IN (1, 2);
	COMMIT;
\end{lstlisting}

\subsection{Pengujian Konkurensi Tinggi}
Untuk menguji bagaimana transaksi berperilaku pada konkurensi tinggi, gunakan script benchmark:

\begin{lstlisting}[language=bash]
	# Di terminal
	cd scripts
	bash benchmark.sh ../praktikum/praktikum2_transaction/transfer_money.sql 20 100 10
\end{lstlisting}

\subsection{Troubleshooting}
Jika mengalami error duplicate key, cek sequence id pada tabel tersebut kemudian reset ke max dari id tabel:

\begin{lstlisting}[language=SQL]
	-- tableName_columnName_seq
	SELECT currval('orders_order_id_seq');
	-- untuk reset
	SELECT setval('orders_order_id_seq', (SELECT MAX(order_id) FROM orders));
\end{lstlisting}

\section{Tugas Praktikum}

\begin{enumerate}
    \item Implementasikan sistem pemesanan produk dengan transaksi yang:
    \begin{itemize}
        \item Mengurangi stok produk
        \item Membuat pesanan baru
        \item Menambahkan item pesanan
        \item Memastikan stok mencukupi
        \item Harus ACID-compliant
    \end{itemize}
    
    \item Implementasikan dua strategi konkurensi berbeda dan bandingkan performanya:
    \begin{itemize}
        \item Row-level locking (pessimistic)
        \item Optimistic concurrency control
    \end{itemize}
    
    \item Lakukan pengujian konkurensi tinggi (minimal 20 koneksi simultan) dan analisis:
    \begin{itemize}
        \item Throughput (transactions per second)
        \item Deadlock yang terjadi (jika ada)
        \item Kesalahan konkurensi lainnya
    \end{itemize}
    
    \item Buat laporan yang menjelaskan bagaimana implementasi transaksi mempengaruhi konkurensi tinggi
\end{enumerate}

\section{Pertanyaan Diskusi}

\begin{enumerate}
    \item Apa perbedaan antara READ COMMITTED dan REPEATABLE READ? Kapan sebaiknya menggunakan masing-masing?
    \item Bagaimana cara mendeteksi dan mengatasi deadlock pada PostgreSQL?
    \item Apa trade-off antara optimistic dan pessimistic concurrency control?
    \item Bagaimana dampak tingkat isolasi transaksi terhadap performa sistem dengan konkurensi tinggi?
    \item Bagaimana strategi yang tepat untuk menangani sistem dengan write-heavy workload?
\end{enumerate}

\section{Referensi}

\begin{itemize}
    \item \href{https://www.postgresql.org/docs/current/transaction-iso.html}{PostgreSQL Documentation: Transactions}
    \item \href{https://www.postgresql.org/docs/current/mvcc.html}{PostgreSQL: Concurrency Control}
    \item \href{https://www.postgresql.org/docs/current/explicit-locking.html}{PostgreSQL: Explicit Locking}
    \item \href{https://www.postgresql.org/docs/current/functions-admin.html}{PostgreSQL: Administrative Functions}
    \item \href{https://www.postgresql.org/docs/current/monitoring-stats.html}{PostgreSQL: Monitoring Database Activity}
\end{itemize}

% Modul P3
\setcounter{section}{0}
\chapter{Implementasi Partitioning pada PostgreSQL}
\setcounter{section}{0}
\section{Pendahuluan}
Partitioning adalah teknik untuk membagi tabel besar menjadi beberapa bagian yang lebih kecil secara fisik, namun tetap terlihat sebagai satu tabel secara logis. Dalam sistem database dengan volume data yang besar, partitioning menjadi strategi penting untuk meningkatkan performa, maintainability, dan skalabilitas.

Ketika volume data tumbuh, performa query dapat menurun signifikan karena database harus memproses lebih banyak data. Dengan partitioning, database dapat membatasi jumlah data yang diakses saat menjalankan query, sehingga meningkatkan efisiensi operasi.

\subsection{Tujuan Praktikum}
Dilaksanakannya praktikum partitioning, praktikan diharapkan mampu:
\begin{enumerate}
    \item Memahami konsep partitioning pada database
    \item Mengimplementasikan berbagai tipe partitioning pada PostgreSQL
    \item Menganalisis dampak partitioning terhadap performa query
    \item Mempelajari strategi partitioning yang tepat untuk berbagai jenis data
\end{enumerate}

\section{Konsep Dasar Partitioning}

\subsection{Apa itu Partitioning?}
Partitioning adalah teknik untuk membagi tabel besar menjadi beberapa bagian yang lebih kecil secara fisik tetapi tetap terlihat sebagai satu tabel secara logis. Hal ini memungkinkan database untuk mengoptimalkan query dengan hanya mengakses partisi yang relevan dengan query, bukan seluruh tabel.

\subsection{Manfaat Partitioning}
\begin{enumerate}
    \item \textbf{Peningkatan Performa Query}: Query yang hanya mengakses subset data (misalnya satu bulan dari data time-series) dapat berjalan lebih cepat karena hanya mengakses partisi yang relevan
    \item \textbf{Efisiensi Maintenance}: Operasi maintenance seperti backup, restore, dan archiving dapat dilakukan pada level partisi
    \item \textbf{Optimalisasi Penyimpanan}: Data dengan pola akses berbeda dapat disimpan pada media penyimpanan yang berbeda (misalnya data terbaru pada SSD)
    \item \textbf{Paralelisme}: Query terhadap beberapa partisi dapat dieksekusi secara paralel
    \item \textbf{Penghapusan Data Efisien}: Menghapus data dalam jumlah besar dapat dilakukan dengan DROP PARTITION daripada DELETE yang lebih lambat
\end{enumerate}

\subsection{Kapan Menggunakan Partitioning}
Partitioning sangat berguna pada:
\begin{enumerate}
    \item Tabel dengan volume data sangat besar (puluhan GB atau lebih)
    \item Data dengan pola akses yang jelas berdasarkan kolom tertentu
    \item Data time-series yang memiliki pola pemasukan dan penghapusan teratur
    \item Tabel fakta pada data warehouse
    \item Sistem yang memerlukan high-availability untuk data terbaru sementara data historis dapat diakses lebih lambat
\end{enumerate}

\subsection{Kapan Tidak Menggunakan Partitioning}
Partitioning tidak selalu bermanfaat pada:
\begin{enumerate}
    \item Tabel dengan ukuran kecil atau menengah yang dapat diproses efisien tanpa partitioning
    \item Tabel yang sering di-join dengan tabel lain (trade-off performa join vs. partitioning)
    \item Sistem dengan pola query yang sulit diprediksi atau tidak selaras dengan strategi partitioning
    \item Ketika overhead administrasi partisi lebih besar dari manfaatnya
\end{enumerate}

\section{Jenis-jenis Partitioning di PostgreSQL}

PostgreSQL menyediakan beberapa tipe partitioning yang dapat dipilih sesuai dengan karakteristik data dan pola akses:

\subsection{Range Partitioning}
\begin{enumerate}
    \item Membagi data berdasarkan rentang nilai pada kolom tertentu
    \item Ideal untuk data time-series (tanggal, timestamp), nilai numerik berurutan
    \item Sangat efektif untuk query yang mengakses data dalam rentang tertentu
    \item Contoh: partisi data penjualan per bulan, per kuartal, atau per tahun
\end{enumerate}

\subsection{List Partitioning}
\begin{enumerate}
    \item Membagi data berdasarkan daftar nilai diskrit pada kolom tertentu
    \item Cocok untuk kolom dengan nilai terbatas dan terdefinisi
    \item Efektif untuk query yang mem-filter berdasarkan nilai spesifik
    \item Contoh: partisi berdasarkan region, status pesanan, kategori produk
\end{enumerate}

\subsection{Hash Partitioning}
\begin{enumerate}
    \item Membagi data berdasarkan hasil fungsi hash dari nilai kolom
    \item Mendistribusikan data secara relatif merata di semua partisi
    \item Cocok saat tidak ada pola akses yang jelas atau untuk distribusi beban
    \item Kurang optimal untuk query range, tetapi baik untuk paralelisme
    \item Contoh: partisi berdasarkan hash dari customer\_id untuk distribusi merata
\end{enumerate}

\begin{figure}[h]
	\centering
	\begin{tabular}{|l|p{10cm}|}
		\hline
		\textbf{Tipe Partitioning} & \textbf{Karakteristik dan Penggunaan} \\
		\hline
		Range & Membagi berdasarkan rentang nilai (tanggal, timestamp, numerik). Ideal untuk data time-series dan query range. Misalnya: data per bulan, per semester, atau per tahun. \\
		\hline
		List & Membagi berdasarkan daftar nilai diskrit. Cocok untuk kolom dengan nilai terbatas seperti region, status, kategori. Misalnya: data per provinsi, per status pesanan. \\
		\hline
		Hash & Membagi berdasarkan hasil fungsi hash. Distribusi data merata, cocok untuk paralelisme dan distribusi beban. Misalnya: hash dari customer\_id untuk distribusi merata. \\
		\hline
	\end{tabular}
	\caption{Jenis-jenis Partitioning di PostgreSQL dan Penggunaannya}
\end{figure}

\subsection{Sub-partitioning (Multi-level Partitioning)}
PostgreSQL juga mendukung partitioning bertingkat (sub-partitioning), di mana partisi dapat dibagi lagi menjadi sub-partisi. Hal ini memungkinkan fleksibilitas lebih tinggi dalam manajemen data:

\begin{enumerate}
    \item Partisi level pertama bisa menggunakan Range (misalnya, per tahun)
    \item Sub-partisi bisa menggunakan List (misalnya, per region atau status)
    \item Atau kombinasi lain dari jenis partitioning
\end{enumerate}

\section{Best Practices Implementasi Partitioning}

\begin{enumerate}
    \item \textbf{Pilih kolom partitioning yang sesuai} dengan pola akses query yang paling umum
    \item \textbf{Batasi jumlah partisi} menjadi jumlah yang dapat dikelola (tidak terlalu banyak)
    \item \textbf{Pertimbangkan partitioning default} untuk menangani data yang tidak masuk kategori partisi yang ada
    \item \textbf{Buat index pada setiap partisi} untuk kolom yang sering digunakan dalam query
    \item \textbf{Gunakan constraint exclusion} untuk memastikan query optimizer melewatkan partisi yang tidak relevan
    \item \textbf{Pertimbangkan tools otomatisasi} untuk mengelola partisi yang dibuat secara berkala
    \item \textbf{Uji performa} untuk memverifikasi bahwa partitioning memberikan manfaat yang diharapkan
    \item \textbf{Kombinasikan dengan strategi penyimpanan} (tablespace) untuk pengoptimalan lebih lanjut
\end{enumerate}

\section{Tahapan Praktikum}

\subsection{Persiapan}
\begin{enumerate}
	\item Pastikan container Docker PostgreSQL sudah berjalan
	\item Gunakan pgAdmin atau PSQL untuk mengakses database
	\item Pastikan data sudah diimpor menggunakan script yang disediakan
\end{enumerate}

\subsection{Implementasi Range Partitioning}
Pada bagian ini, kita akan membuat partisi tabel orders berdasarkan tanggal pesanan (order date):

\begin{lstlisting}[language=SQL]
	-- Buat tabel induk dengan deklarasi partisi
	CREATE TABLE orders_partitioned (
		order_id SERIAL,
		customer_id INTEGER NOT NULL,
		order_date TIMESTAMP NOT NULL,
		status VARCHAR(50) DEFAULT 'pending',
		total_amount DECIMAL(12, 2),
		shipping_address TEXT,
		payment_method VARCHAR(50),
		PRIMARY KEY (order_id, order_date),
		CONSTRAINT valid_status CHECK (status IN ('pending', 'processing', 'shipped', 'delivered', 'cancelled'))
	) PARTITION BY RANGE (order_date);
	
	-- Buat partisi untuk setiap kuartal tahun 2023
	CREATE TABLE orders_q1_2023 PARTITION OF orders_partitioned
		FOR VALUES FROM ('2023-01-01') TO ('2023-04-01');
	
	CREATE TABLE orders_q2_2023 PARTITION OF orders_partitioned
		FOR VALUES FROM ('2023-04-01') TO ('2023-07-01');
	
	CREATE TABLE orders_q3_2023 PARTITION OF orders_partitioned
		FOR VALUES FROM ('2023-07-01') TO ('2023-10-01');
	
	CREATE TABLE orders_q4_2023 PARTITION OF orders_partitioned
		FOR VALUES FROM ('2023-10-01') TO ('2024-01-01');
	
	-- Import data dari tabel orders asli
	INSERT INTO orders_partitioned
	SELECT * FROM orders;
\end{lstlisting}


\subsection{Implementasi List Partitioning}
Selanjutnya, kita akan mencoba partitioning berdasarkan status pesanan:

\begin{lstlisting}[language=SQL]
	-- Buat tabel induk dengan deklarasi partisi
	CREATE TABLE orders_by_status (
		order_id SERIAL,
		customer_id INTEGER NOT NULL,
		order_date TIMESTAMP NOT NULL,
		status VARCHAR(50) NOT NULL,
		total_amount DECIMAL(12, 2),
		shipping_address TEXT,
		payment_method VARCHAR(50),
		PRIMARY KEY (order_id, status)
	) PARTITION BY LIST (status);
	
	-- Buat partisi untuk setiap status
	CREATE TABLE orders_pending PARTITION OF orders_by_status
		FOR VALUES IN ('pending');
	
	CREATE TABLE orders_processing PARTITION OF orders_by_status
		FOR VALUES IN ('processing');
	
	CREATE TABLE orders_shipped PARTITION OF orders_by_status
		FOR VALUES IN ('shipped');
	
	CREATE TABLE orders_delivered PARTITION OF orders_by_status
		FOR VALUES IN ('delivered');
	
	CREATE TABLE orders_cancelled PARTITION OF orders_by_status
		FOR VALUES IN ('cancelled');
	
	-- Import data dari tabel orders asli
	INSERT INTO orders_by_status
	SELECT * FROM orders;
\end{lstlisting}

\subsection{Implementasi Hash Partitioning}
Berikutnya, kita akan mencoba partitioning berdasarkan hash dari customer id:

\begin{lstlisting}[language=SQL]
	-- Buat tabel induk dengan deklarasi partisi
	CREATE TABLE orders_by_customer (
		order_id SERIAL,
		customer_id INTEGER NOT NULL,
		order_date TIMESTAMP NOT NULL,
		status VARCHAR(50) DEFAULT 'pending',
		total_amount DECIMAL(12, 2),
		shipping_address TEXT,
		payment_method VARCHAR(50),
		PRIMARY KEY (order_id, customer_id),
		CONSTRAINT valid_status CHECK (status IN ('pending', 'processing', 'shipped', 'delivered', 'cancelled'))
	) PARTITION BY HASH (customer_id);
	
	-- Buat 4 partisi
	CREATE TABLE orders_by_customer_0 PARTITION OF orders_by_customer
		FOR VALUES WITH (MODULUS 4, REMAINDER 0);
	
	CREATE TABLE orders_by_customer_1 PARTITION OF orders_by_customer
		FOR VALUES WITH (MODULUS 4, REMAINDER 1);
	
	CREATE TABLE orders_by_customer_2 PARTITION OF orders_by_customer
		FOR VALUES WITH (MODULUS 4, REMAINDER 2);
	
	CREATE TABLE orders_by_customer_3 PARTITION OF orders_by_customer
		FOR VALUES WITH (MODULUS 4, REMAINDER 3);
	
	-- Import data dari tabel orders asli
	INSERT INTO orders_by_customer
	SELECT * FROM orders;
\end{lstlisting}

\subsection{Menganalisis Performa Partitioning}
Pada bagian ini, kita akan membandingkan performa query antara tabel biasa dan tabel yang menggunakan partitioning:

\begin{lstlisting}[language=SQL]
	-- Query pada tabel non-partisi
	EXPLAIN ANALYZE
	SELECT *
	FROM orders
	WHERE order_date BETWEEN '2023-06-01' AND '2023-06-30';
	
	-- Query pada tabel dengan partisi range
	EXPLAIN ANALYZE
	SELECT *
	FROM orders_partitioned
	WHERE order_date BETWEEN '2023-06-01' AND '2023-06-30';
\end{lstlisting}

\subsection{Maintenance Partisi}
Selanjutnya, kita akan mempelajari bagaimana melakukan maintenance terhadap partisi:

\begin{lstlisting}[language=SQL]
	-- Menambahkan partisi baru untuk tahun 2024
	CREATE TABLE orders_q1_2024 PARTITION OF orders_partitioned
		FOR VALUES FROM ('2024-01-01') TO ('2024-04-01');
	
	-- Menghapus data lama dari partisi tertentu
	TRUNCATE TABLE orders_q1_2023;
	
	-- Mengarsipkan partisi lama
	ALTER TABLE orders_partitioned DETACH PARTITION orders_q1_2023;
	-- Sekarang orders_q1_2023 adalah tabel mandiri yang dapat diarchive
\end{lstlisting}

\subsection{Implementasi Sub-partitioning (Multi-level Partitioning)}
Kita juga akan mengimplementasikan partitioning bertingkat:

\begin{lstlisting}[language=SQL]
	-- Buat tabel induk dengan deklarasi partisi bertingkat
	CREATE TABLE orders_multi_level (
		order_id SERIAL,
		customer_id INTEGER NOT NULL,
		order_date TIMESTAMP NOT NULL,
		status VARCHAR(50) NOT NULL,
		total_amount DECIMAL(12, 2),
		shipping_address TEXT,
		payment_method VARCHAR(50),
		PRIMARY KEY (order_id, order_date, status)
	) PARTITION BY RANGE (order_date);
	
	-- Buat partisi kuartal, yang juga akan dipartisi lagi
	CREATE TABLE orders_q1_2023_multi PARTITION OF orders_multi_level
		FOR VALUES FROM ('2023-01-01') TO ('2023-04-01')
		PARTITION BY LIST (status);
	
	-- Subpartisi berdasarkan status untuk Q1 2023
	CREATE TABLE orders_q1_2023_pending PARTITION OF orders_q1_2023_multi
		FOR VALUES IN ('pending');
	
	CREATE TABLE orders_q1_2023_processing PARTITION OF orders_q1_2023_multi
		FOR VALUES IN ('processing');
	
	CREATE TABLE orders_q1_2023_shipped PARTITION OF orders_q1_2023_multi
		FOR VALUES IN ('shipped');
	
	CREATE TABLE orders_q1_2023_delivered PARTITION OF orders_q1_2023_multi
		FOR VALUES IN ('delivered');
	
	CREATE TABLE orders_q1_2023_cancelled PARTITION OF orders_q1_2023_multi
		FOR VALUES IN ('cancelled');
	
	-- Import data
	INSERT INTO orders_multi_level
	SELECT * FROM orders 
	WHERE order_date BETWEEN '2023-01-01' AND '2023-03-31';
\end{lstlisting}

\section{Tugas Praktikum}

\begin{enumerate}
    \item Implementasikan partitioning pada salah satu tabel besar (orders atau products)
    
    \item Bandingkan performa query antara tabel non-partisi dan tabel dengan partisi untuk:
    \begin{itemize}
        \item Query yang mengakses satu partisi
        \item Query yang mengakses beberapa partisi
        \item Query yang mengakses seluruh tabel
    \end{itemize}
    
    \item Implementasikan minimal dua tipe partitioning berbeda (Range, List, atau Hash)
    
    \item Buat script untuk maintenance partisi secara otomatis:
    \begin{itemize}
        \item Menambahkan partisi baru secara otomatis
        \item Mengarsipkan partisi lama
    \end{itemize}
    
    \item Lakukan benchmark konkuren dan bandingkan throughput pada tabel dengan dan tanpa partisi
\end{enumerate}

\section{Pertanyaan Diskusi}

\begin{enumerate}
    \item Kapan sebaiknya menggunakan partitioning dan kapan tidak perlu?
    
    \item Bagaimana strategi partitioning yang tepat untuk:
    \begin{itemize}
        \item Data time-series dengan volume tinggi
        \item Data yang sering diakses berdasarkan wilayah geografis
        \item Tabel fakta pada data warehouse
    \end{itemize}
    
    \item Apa perbedaan antara partitioning dan sharding? Kapan sebaiknya menggunakan masing-masing?
    
    \item Apa kelebihan dan kekurangan dari masing-masing tipe partitioning (Range, List, Hash)?
    
    \item Bagaimana partitioning mempengaruhi performa JOIN operations?
\end{enumerate}

\section{Referensi}

\begin{itemize}
    \item \href{https://www.postgresql.org/docs/current/ddl-partitioning.html}{PostgreSQL Documentation: Table Partitioning}
    \item \href{https://www.postgresql.org/docs/current/ddl-partitioning.html#DDL-PARTITIONING-DECLARATIVE}{PostgreSQL: Declarative Partitioning}
    \item \href{https://wiki.postgresql.org/wiki/Table_partitioning}{PostgreSQL Wiki: Table Partitioning}
    \item \href{https://www.postgresql.org/docs/current/ddl-partitioning.html#DDL-PARTITIONING-IMPLEMENTATION}{PostgreSQL: Partition Pruning}
    \item \href{https://www.postgresql.org/docs/current/ddl-partitioning.html#DDL-PARTITIONING-CONSTRAINT-EXCLUSION}{PostgreSQL: Constraint Exclusion}
\end{itemize}

% Modul P4
\setcounter{section}{0}
\chapter{Backup dan Recovery pada PostgreSQL}
\setcounter{section}{0}
\section{Pendahuluan}
Backup dan recovery merupakan aspek kritis dalam pengelolaan database. Sebagai database administrator, melindungi data dari kehilangan dan kerusakan merupakan tanggung jawab utama. PostgreSQL menyediakan beragam metode backup dan recovery yang dapat disesuaikan dengan kebutuhan sistem, batasan operasional, dan tujuan bisnis.

Strategi backup dan recovery yang tepat menjamin ketahanan (resilience) dan ketersediaan (availability) sistem database. Dalam lingkungan produksi, kemampuan untuk memulihkan data dengan cepat dan akurat saat terjadi kegagalan sistem, kesalahan manusia, atau bencana dapat menjadi perbedaan antara operasi bisnis yang berkelanjutan dan kerugian finansial yang signifikan.

\subsection{Tujuan Praktikum}
Dilaksanakannya praktikum backup dan recovery, praktikan diharapkan mampu:
\begin{enumerate}
    \item Memahami konsep backup dan recovery pada database
    \item Mengimplementasikan berbagai tipe backup di PostgreSQL
    \item Melakukan recovery data dari backup
    \item Menerapkan strategi backup yang tepat sesuai kebutuhan
\end{enumerate}

\section{Konsep Dasar Backup dan Recovery}

\subsection{Terminologi dan Konsep Kunci}
\begin{enumerate}
    \item \textbf{Backup}: Proses membuat salinan data untuk pemulihan jika terjadi kehilangan data
    \item \textbf{Recovery}: Proses memulihkan database ke keadaan yang konsisten menggunakan data backup
    \item \textbf{Recovery Point Objective (RPO)}: Jumlah maksimal data yang dapat hilang, diukur dalam waktu
    \item \textbf{Recovery Time Objective (RTO)}: Waktu maksimal yang diperbolehkan untuk memulihkan sistem
    \item \textbf{Full Backup}: Backup seluruh database
    \item \textbf{Incremental Backup}: Backup hanya perubahan sejak backup terakhir
    \item \textbf{Differential Backup}: Backup perubahan sejak full backup terakhir
    \item \textbf{Point-in-Time Recovery (PITR)}: Kemampuan memulihkan ke titik waktu tertentu
\end{enumerate}

\subsection{Tipe-tipe Backup di PostgreSQL}
PostgreSQL menyediakan beberapa metode backup yang dapat digunakan sesuai kebutuhan:

\begin{enumerate}
    \item \textbf{Logical Backup}: 
    \begin{itemize}
        \item Menggunakan utilitas \texttt{pg\_dump} dan \texttt{pg\_dumpall}
        \item Menghasilkan file SQL atau format khusus yang berisi perintah untuk merekonstruksi database
        \item Fleksibel dan dapat dipulihkan secara selektif
        \item Cocok untuk database kecil hingga menengah
    \end{itemize}
    
    \item \textbf{Physical Backup}:
    \begin{itemize}
        \item Salinan langsung dari file data PostgreSQL
        \item Lebih cepat untuk database besar
        \item Membutuhkan akses ke sistem file
        \item Membutuhkan penghentian database atau konsistensi filesystem
    \end{itemize}
    
    \item \textbf{Continuous Archiving}:
    \begin{itemize}
        \item Menggunakan Write-Ahead Logs (WAL) untuk mencatat semua perubahan
        \item Memungkinkan Point-in-Time Recovery (PITR)
        \item Kombinasi dari base backup dan WAL archives
        \item Cocok untuk database dengan zero/minimal data loss requirement
    \end{itemize}
\end{enumerate}

\begin{figure}[h]
	\centering
	\begin{tabular}{|l|p{4cm}|p{4cm}|p{4cm}|}
		\hline
		\textbf{Aspek} & \textbf{Logical Backup} & \textbf{Physical Backup} & \textbf{Continuous Archiving} \\
		\hline
		Metode & pg\_dump, pg\_dumpall & File system backup & Base backup + WAL archiving \\
		\hline
		Kelebihan & Selektif, platform-independent, dapat restore sebagian & Cepat untuk database besar, overhead rendah & Point-in-time recovery, minimal data loss \\
		\hline
		Kekurangan & Lambat untuk database besar, overhead tinggi & Tidak selektif, platform-dependent & Kompleksitas setup, kebutuhan penyimpanan tinggi \\
		\hline
		Ideal untuk & Database kecil-menengah, backup selektif & Database besar, full restore cepat & Sistem kritikal, zero data loss requirement \\
		\hline
	\end{tabular}
	\caption{Perbandingan Tipe Backup di PostgreSQL}
\end{figure}

\section{Strategi Backup yang Efektif}

\subsection{Faktor yang Mempengaruhi Strategi Backup}
\begin{enumerate}
    \item \textbf{Volume Data}: Ukuran database mempengaruhi waktu backup dan storage requirements
    \item \textbf{Window Backup}: Periode waktu yang tersedia untuk melakukan backup
    \item \textbf{RPO dan RTO}: Persyaratan bisnis untuk toleransi kehilangan data dan waktu pemulihan
    \item \textbf{Resource Constraints}: Keterbatasan penyimpanan, CPU, dan bandwith jaringan
    \item \textbf{Frekuensi Perubahan}: Seberapa sering dan banyak data berubah
    \item \textbf{Retensi Backup}: Berapa lama backup harus disimpan
\end{enumerate}

\subsection{Rekomendasi Best Practices}
\begin{enumerate}
    \item \textbf{Redundansi}: Simpan beberapa salinan backup di lokasi berbeda
    \item \textbf{Verifikasi Backup}: Secara teratur menguji integritas backup dengan test restore
    \item \textbf{Otomatisasi}: Jadwalkan backup secara reguler menggunakan cron jobs atau tools
    \item \textbf{Monitoring}: Pantau proses backup dan segera tangani kegagalan
    \item \textbf{Dokumentasi}: Dokumentasikan prosedur backup dan recovery secara detail
    \item \textbf{Enkripsi}: Lindungi data backup dengan enkripsi jika berisi informasi sensitif
    \item \textbf{Rotasi Backup}: Implementasikan strategi rotasi untuk mengoptimalkan penyimpanan
    \item \textbf{Strategi Berlapis}: Kombinasikan beberapa metode backup (misal: full backup mingguan + incremental harian + continuous archiving)
\end{enumerate}

\section{Tahapan Praktikum}

\subsection{Persiapan}
\begin{enumerate}
	\item Pastikan container Docker PostgreSQL sudah berjalan
	\item Gunakan pgAdmin atau PSQL untuk mengakses database
	\item Pastikan data sudah diimpor menggunakan script yang disediakan
\end{enumerate}

\subsection{Logical Backup dengan pg\_dump}
Mari kita mulai dengan melakukan logical backup menggunakan pg\_dump:

\subsubsection{Backup Database Lengkap}
\begin{lstlisting}[language=bash]
	# Format SQL
	pg_dump -h localhost -p 5432 -U praktikan -d ecommerce -f ecommerce_full.sql
	
	# Format Custom (untuk restore selektif)
	pg_dump -h localhost -p 5432 -U praktikan -d ecommerce -Fc -f ecommerce_full.dump
	
	# Format Directory (paralel backup)
	pg_dump -h localhost -p 5432 -U praktikan -d ecommerce -Fd -j 4 -f ecommerce_backup_dir
\end{lstlisting}

\subsubsection{Backup Tabel Tertentu}
\begin{lstlisting}[language=bash]
	# Backup hanya tabel orders dan order_items
	pg_dump -h localhost -p 5432 -U praktikan -d ecommerce -t orders -t order_items -f orders_backup.sql
\end{lstlisting}

\subsubsection{Backup Schema Saja}
\begin{lstlisting}[language=bash]
	# Backup struktur tanpa data
	pg_dump -h localhost -p 5432 -U praktikan -d ecommerce --schema-only -f ecommerce_schema.sql
\end{lstlisting}

\subsubsection{Backup Data Saja}
\begin{lstlisting}[language=bash]
	# Backup data tanpa struktur
	pg_dump -h localhost -p 5432 -U praktikan -d ecommerce --data-only -f ecommerce_data.sql
\end{lstlisting}

\subsection{Restore dari Logical Backup}
Setelah melakukan backup, kita juga perlu tahu cara melakukan restore:

\subsubsection{Restore Database Lengkap}
\begin{lstlisting}[language=bash]
	# Restore dari format SQL
	psql -h localhost -p 5432 -U praktikan -d ecommerce_restore -f ecommerce_full.sql
	
	# Restore dari format Custom
	pg_restore -h localhost -p 5432 -U praktikan -d ecommerce_restore -v ecommerce_full.dump
\end{lstlisting}

\subsubsection{Restore Tabel Tertentu}
\begin{lstlisting}[language=bash]
	# Restore hanya tabel orders dari backup custom
	pg_restore -h localhost -p 5432 -U praktikan -d ecommerce_restore -t orders -v ecommerce_full.dump
\end{lstlisting}

\subsection{Physical Backup}
Physical backup melibatkan penyalinan langsung file data PostgreSQL. Dalam container Docker, kita bisa menggunakan volume untuk menyimpan data:

\begin{lstlisting}[language=bash]
	# Stop container
	docker-compose stop postgres
	
	# Backup volume
	docker run --rm -v postgres_praktikum_postgres_data:/source -v $(pwd)/backup:/backup ubuntu tar -czf /backup/postgres_data_backup.tar.gz -C /source .
	
	# Start container kembali
	docker-compose start postgres
\end{lstlisting}

\subsection{Continuous Archiving dan Point-in-Time Recovery (PITR)}
Untuk mengimplementasikan Continuous Archiving, kita perlu mengkonfigurasi PostgreSQL untuk menyimpan Write-Ahead Logs (WAL):

\subsubsection{Konfigurasi WAL Archiving}
Edit file \texttt{postgresql.conf}:

\begin{lstlisting}
	wal_level = replica
	archive_mode = on
	archive_command = 'cp %p /var/lib/postgresql/data/archive/%f'
\end{lstlisting}

\subsubsection{Point-in-Time Recovery}
\begin{lstlisting}[language=bash]
	# Buat base backup
	pg_basebackup -h localhost -p 5432 -U praktikan -D /backup/base -Fp -Xs -P
	
	# Restore ke titik waktu tertentu
	# (Memerlukan recovery.conf)
\end{lstlisting}

\subsection{Demonstrasi Skenario Disaster}
Pada bagian ini, kita akan mensimulasikan skenario kehilangan data dan melakukan recovery:

\subsubsection{Simulasi Kehilangan Data}
\begin{lstlisting}[language=SQL]
	-- Buat tabel untuk demonstrasi
	CREATE TABLE important_data (
		id SERIAL PRIMARY KEY,
		description TEXT,
		created_at TIMESTAMP DEFAULT CURRENT_TIMESTAMP
	);
	
	-- Insert data penting
	INSERT INTO important_data (description) VALUES ('Data sangat penting #1');
	INSERT INTO important_data (description) VALUES ('Data sangat penting #2');
	INSERT INTO important_data (description) VALUES ('Data sangat penting #3');
	
	-- Backup
	-- (Jalankan pg_dump seperti di atas)
	
	-- Simulasi kesalahan
	DROP TABLE important_data;
	
	-- Restore
	-- (Jalankan pg_restore seperti di atas)
\end{lstlisting}

\subsubsection{Simulasi Corrupt Database}
\begin{lstlisting}[language=bash]
	# Backup terlebih dahulu
	# (Jalankan backup)
	
	# Simulasi korupsi (hati-hati!)
	docker-compose exec postgres bash -c "dd if=/dev/urandom of=/var/lib/postgresql/data/base/16384/1234 bs=8192 count=1 conv=notrunc"
	
	# Restart PostgreSQL (kemungkinan akan gagal start)
	docker-compose restart postgres
	
	# Restore dari backup
	# (Jalankan restore)
\end{lstlisting}

\subsection{Strategi Backup yang Baik}
Berdasarkan praktik terbaik, berikut adalah strategi backup yang direkomendasikan:
\begin{enumerate}
	\item Full backup setiap minggu
	\item Differential backup setiap hari
	\item Continuous archiving untuk recovery point-in-time
	\item Replikasi untuk high availability
\end{enumerate}

\subsection{Implementasi High Availability dengan Replikasi}
Meskipun bukan bagian dari backup/recovery secara tradisional, replikasi membantu memastikan ketersediaan data:

\begin{lstlisting}
	# Primary server (postgresql.conf)
	wal_level = replica
	max_wal_senders = 10
	wal_keep_segments = 64
	
	# Standby server (postgresql.conf)
	hot_standby = on
	
	# Standby server (recovery.conf)
	standby_mode = 'on'
	primary_conninfo = 'host=primary port=5432 user=replicator password=password'
	trigger_file = '/tmp/postgresql.trigger'
\end{lstlisting}

\section{Tugas Praktikum}

\begin{enumerate}
    \item Lakukan backup database dengan minimal 3 metode berbeda:
    \begin{itemize}
        \item Full database backup dengan pg\_dump
        \item Backup tabel tertentu
        \item Backup struktur saja
    \end{itemize}
    
    \item Simulasikan skenario kehilangan data dan lakukan recovery dari backup:
    \begin{itemize}
        \item Menghapus tabel penting secara tidak sengaja
        \item Menghapus beberapa baris data
        \item Menghapus database
    \end{itemize}
    
    \item Buat script otomatisasi backup yang berjalan terjadwal:
    \begin{itemize}
        \item Menggunakan cron job atau scheduled tasks
        \item Menyimpan backup dengan timestamp
        \item Menghapus backup lama secara otomatis
    \end{itemize}
    
    \item Buat laporan yang menjelaskan strategi backup dan recovery yang tepat untuk database dengan karakteristik:
    \begin{itemize}
        \item Data transaksi finansial yang kritikal
        \item Data logging dengan volume tinggi
        \item Data referensi yang jarang berubah
    \end{itemize}
\end{enumerate}

\section{Pertanyaan Diskusi}

\begin{enumerate}
    \item Apa kelebihan dan kekurangan masing-masing metode backup (logical vs physical vs continuous archiving)?
    
    \item Bagaimana strategi backup yang tepat untuk database dengan ukuran sangat besar (>1TB)?
    
    \item Apa trade-off antara keamanan data (durability) dengan performa database?
    
    \item Bagaimana cara menghitung Recovery Point Objective (RPO) dan Recovery Time Objective (RTO) untuk database?
    
    \item Apa saja pertimbangan saat merancang strategi disaster recovery untuk database produksi?
\end{enumerate}

\section{Referensi}

\begin{itemize}
    \item \href{https://www.postgresql.org/docs/current/backup.html}{PostgreSQL Documentation: Backup and Restore}
    \item \href{https://www.postgresql.org/docs/current/continuous-archiving.html}{PostgreSQL: Continuous Archiving and PITR}
    \item \href{https://www.postgresql.org/docs/current/high-availability.html}{PostgreSQL: High Availability}
    \item \href{https://wiki.postgresql.org/wiki/Backup_in_practice}{PostgreSQL Wiki: Backup in Practice}
    \item \href{https://www.postgresql.org/docs/current/app-pgdump.html}{PostgreSQL: pg\_dump}
    \item \href{https://www.postgresql.org/docs/current/app-pgrestore.html}{PostgreSQL: pg\_restore}
\end{itemize}

\end{document}